\section{Proofs for Section~\ref{sec:soft}}\label{app:soft}
Notation and the labels ``proved'', ``by computer'' and ``evidence'' are as in Section~\ref{sec:soft}. Computer checks were carried out by two independent implementations, the second of which built every instance and every set $\Sigma$ directly from the graph and the definition (all set partitions, or an exact subset dynamic program), without using the structure lemmas.

\subsection{Proof of Theorem~\ref{thm:sfA}}\label{app:softA}
By Lemma~\ref{lem:sfstruct}, $\Sigma$ is empty or the intersection of unary lists and of binary relations $R_{AB}$; it suffices to show that each of these is closed under every majority operation $m$. A static violation makes $\Sigma$ empty, which is closed.

\emph{Lists.} Let $S\subseteq\{1,2,3\}$ and $u,v,w\in S$. If two of them coincide then $m(u,v,w)$ is their common value, which lies in $S$. If all three are distinct then $S=\{1,2,3\}$.

\emph{Binary relations (classification by computer).} $R_{AB}$ is determined by the roles of the shared vertex $q$ in $A$ and in $B$ (three each) and the four cross bits, $144$ configurations. Computing the relation from the formula of Lemma~\ref{lem:sfstruct}(iii) gives exactly $40$ distinct relations, with the following shapes, where $p,q$ range over $\{1,2\}$ (the states that edit a pair at the shared vertex):
\begin{center}
\begin{tabular}{@{}lc@{}}
shape & number\\\hline
$\emptyset$ & $1$\\
singletons & $9$\\
two-element sets & $18$\\
rows $\{p\}\times\{1,2,3\}$, columns $\{1,2,3\}\times\{q\}$ ($4$) and the $3$-element subsets of $\{1,2\}^2$ ($4$) & $8$\\
crosses $\bigl(\{p\}\times\{1,2,3\}\bigr)\cup\bigl(\{1,2,3\}\times\{q\}\bigr)$ & $4$
\end{tabular}
\end{center}
The same list results from an independent construction: all $576$ graphs on five vertices consisting of two members sharing exactly one vertex (all $6\times6$ orientations, $16$ patterns of the four cross pairs), with each triple $\{p,q,r\}$ tested for being an induced $P_3$ after editing, without the formula.

\emph{Closure of each shape.} Let $m$ be a majority and $t_1,t_2,t_3\in R$, $u=m(t_1,t_2,t_3)$.
(1) $|R|\le2$: at every coordinate at most two values occur among $t_1,t_2,t_3$, so $u$ is the coordinatewise majority; if, say, $t_1$ occurs at least twice among the three tuples, then $u=t_1\in R$.
(2) A row $\{p\}\times\{1,2,3\}$: the first coordinate is constantly $p$, so $u_1=p$ and $u\in R$; columns are symmetric.
(3) A $3$-element subset of $\{1,2\}^2$: all values lie in $\{1,2\}$, so $m$ acts as Boolean majority, and every binary relation on $\{1,2\}$ is closed under it (it is definable by clauses with at most two literals).
(4) A cross $R=(\{p\}\times\{1,2,3\})\cup(\{1,2,3\}\times\{q\})$: suppose $u_1\ne p$; we show $u_2=q$. If two of the three first coordinates coincide, $u_1$ is their value, which differs from $p$, so these two tuples are not in the row and are of the form $(x,q)$; two second coordinates equal $q$, so $u_2=q$. If the three first coordinates are distinct, at most one of them equals $p$, so at least two tuples are of the form $(x,q)$ and again $u_2=q$. The case $u_2\ne q$ is symmetric. So $u_1=p$ or $u_2=q$, i.e.\ $u\in R$. \hfill$\square$

\emph{Orientation.} The state names $1,2$ depend on the orientation $(a,b,c)$ versus $(c,b,a)$ of a member (they are exchanged). The set of all majority operations is invariant under relabelling of the states by any permutation, so closure under \emph{all} majorities does not depend on this convention (closure under one fixed operation would).

\subsection{Proof of Theorem~\ref{thm:U}}\label{app:softU}
Call a pair of vertices a \emph{cross pair} or \emph{far} as in Appendix~\ref{sec:softfar}. For a graph $G'=G\sd W$ and a member $A=(a,b,c)$ and vertex $t\notin V(A)$ whose three pairs $\{x,t\}$, $x\in V(A)$, are hard, the \emph{list of $(A,t)$} is the set of states of Observation~\ref{obs:lists}. It is \emph{liftable} if at least one of the three pairs is far. A far pair lies in no binary triple (the third pair of a binary triple is a cross pair) and in no triple with two soft pairs. A triple $\{x,y,t\}$ with one soft pair $xy$ of $A=(x,y,z)$ (up to naming) and a far pair $xt$ also has $yt$ and $zt$ hard: if $yt$ (or $zt$) were soft, with owner $B$, then $B$ shares the vertex $y$ (or $z$) with $A$ and $xt$ would be a cross pair. So the triple belongs to the list of $(A,t)$, which is liftable. By Lemma~\ref{lem:sfU} a far pair therefore changes only static triples and liftable lists. Let $\Sigma_{\mathrm{perm}}(G')$ be the set of vectors satisfying all constraints of $G'$ except the liftable lists; it is closed under every majority operation by Theorem~\ref{thm:sfA} (only unary lists and the relations $R_{AB}$ remain), and $\Sigma_{\mathrm{perm}}(G\sd W)=\Sigma_{\mathrm{perm}}(G)$ for every set $W$ of far pairs.

\emph{Algorithm} $\mathrm{solve}(G',b)$, deciding whether some set $W$ of at most $b$ far pairs has $\Sigma(G'\sd W)\ne\emptyset$.
(1) If $G'$ has an induced $P_3$ with three hard pairs, branch on its far pairs, $b\to b-1$ (if there is none, or $b=0$, answer NO).
(2) Else if $\Sigma(G')\ne\emptyset$, answer YES; if $b=0$, NO.
(3) If $\Sigma_{\mathrm{perm}}(G')=\emptyset$, NO.
(4) Else choose a minimal family $\mathcal F$ of liftable lists with $\Sigma_{\mathrm{perm}}(G')\cap\mathcal F=\emptyset$ (greedy deletion). By Lemma~\ref{lem:sfH2}, $|\mathcal F|\le2$ (a list is a unit list, or empty; empty is a family of one). For each list of $\mathcal F$ and each of its far pairs $e$ (at most $3$) call $\mathrm{solve}(G'\sd e,b-1)$.

\emph{Correctness.} Soundness is clear. Completeness by induction on $b$: let $W$ be a smallest set of at most $b$ far pairs with $\Sigma(G'\sd W)\ne\emptyset$. If $W=\emptyset$ step (2) answers YES (static violations are absent by (1) and then $\Sigma\ne\emptyset$ is tested). Otherwise $b\ge1$. A static all-hard induced $P_3$ of $G'$ has to be hit by a far pair of $W$ (only far pairs are edited), so (1) branches correctly. In (4), $W$ changes no relation and no permanent list, so $\Sigma(G'\sd W)\subseteq\Sigma_{\mathrm{perm}}(G')\cap\{\text{unchanged liftable lists}\}$; as $\Sigma(G'\sd W)\ne\emptyset$ and $\Sigma_{\mathrm{perm}}(G')\cap\mathcal F=\emptyset$, $W$ changes some list of $\mathcal F$, i.e.\ contains a far pair $\{x,t\}$ with $x\in V(A)$ of it. Recursing on $G'\sd e$ with $W\setminus\{e\}$ completes the induction (the pair $e$ stays far in $G'\sd e$, and a far pair is a far pair of $H$ whichever graph is considered).

\emph{Time.} Each node takes polynomial time (Corollary~\ref{cor:sfpoly}), at most $6$ children, depth at most $\ell$. \hfill$\square$

\subsection{Algorithm \texorpdfstring{$\mathcal A$}{A} and the proof of Theorem~\ref{thm:F}}\label{app:softF}
Let $I(H)$ be a forest. Write $\mathrm{Ex}_C(G')$ as in Appendix~\ref{sec:softmodel}. For a graph $G'$ (obtained from the input by flipping hard pairs) without an all-hard induced $P_3$, every triple whose three pairs are far is not an induced $P_3$, so it can be deleted without changing $\mathrm{Ex}$; hence $\mathrm{Ex}_C(G')=\mathrm{Ex}(G')$ if $C$ is the set of all sites.

\emph{Algorithm} $\mathcal A(G',r)$ (excess budget $r$; accept iff $\OPT(G')\le|H|+r$).
(1) If $G'$ has an induced $P_3$ with three hard pairs: if $r=0$ reject; else call $\mathcal A(G'\sd e,r-1)$ for each of its three pairs $e$ (cross pairs included; the recursion is on graphs, so this stays sound) and accept iff one call accepts.
(2) Now $H$ is pair-maximal in $G'$. If $\mathrm{Ex}(G')\le r$, accept; if $r=0$, reject.
(3) Compute an inclusion-minimal set $C$ of sites with $\mathrm{Ex}_C(G')>r-1$ by greedy deletion (at most $2|H|$ oracle calls).
(4) With $B=2D+3r$ let $L_B$ be the set of Appendix~\ref{sec:optimal}, computed from $G',H$, and $P(x)$ the set of $y\in\mathrm{Ball}_{G'[L_B]}(x,2)$ with $xy$ far. For every $x\in V(C)$ and $y\in P(x)$ call $\mathcal A(G'\sd xy,r-1)$; accept iff one accepts.

\emph{Soundness.} Acceptance in (2) gives a solution of size $\le|H|+r$ by Lemma~\ref{lem:sfE} (take $W=\emptyset$); in (1) and (4) a solution of $G'\sd e$ of size $\le|H|+r-1$ gives one of $G'$ of size $\le|H|+r$.

\emph{Completeness.} Suppose $\OPT(G')\le|H|+r$, there is no static violation and $\mathrm{Ex}(G')>r$ (otherwise the algorithm accepts). In step (1), a solution contains one of the three pairs of the static $P_3$, and removing it gives a solution of $G'\sd e$ of size $\le|H|+r-1$. In (4): by Lemma~\ref{lem:sfE} choose a set $W$ of far pairs minimising $|W|+\mathrm{Ex}(G'\sd W)\le r$, with minimising $F,x$; $(W,F,x)$ is an optimum of $G'$. As $\mathrm{Ex}(G')>r$, $W\ne\emptyset$ and $\mathrm{Ex}(G'\sd W)\le r-1$. Let $Z(W)$ be the set of members through an end of a pair of $W$ together with all edges of $I(H)$ at these members. Every triple that contains a pair of $W$ has all its sites in $Z(W)$ (an owner of a soft pair of it is a member through an end of the pair; a cross pair $\{x,z\}$ of it lies at an edge incident to a member through $x$ or $z$). Let $C'$ be the set of sites outside $Z(W)$. The triples kept in $\mathrm{Ex}_{C'}(G')$ contain no pair of $W$, so they have the same status in $G'$ and in $G'\sd W$, and $(F,x)$ is feasible for $\mathrm{Ex}_{C'}(G')$ with the same cost: $\mathrm{Ex}_{C'}(G')\le\mathrm{Ex}(G'\sd W)\le r-1$. The set $C$ of step (3) has $\mathrm{Ex}_C(G')>r-1$, so by monotonicity $C\not\subseteq C'$: $C$ contains a site of $Z(W)$, a member through an end $x$ of some $e=xy\in W$, or an edge incident to such a member, and in both cases $x\in V(C)$. Since $(W,F,x)$ is an optimum of $G'$ of size $\le|H|+r$ and $H$ is pair-maximal in $G'$, Corollary~\ref{cor:cand} gives $y\in\mathrm{Ball}(x,2)$ of $G'[L_B]$, so $y\in P(x)$. The recursion on $G'\sd e$ with $W\setminus\{e\}$ satisfies $|W\setminus\{e\}|+\mathrm{Ex}\le r-1$.

\emph{Time.} The depth is at most $r\le\ell$; a node has at most $\max\{3,|V(C)|Z^2\}\le6|C|Z^2$ children, $Z=3D+3r+1$; each node makes at most $2|H|+1$ oracle calls. Correctness did not use BC; only $|C|\le K$ does. \hfill$\square$

\subsection{Proof of Theorem~\ref{thm:L}}\label{app:softL}
Variables of the CSP: the state of each member (seven non-empty edit sets, cost $=$ size) and a bit for each local hard pair (cost $1$ if edited); remote pairs are constants (in $T$) or absent. A vertex triple $\tau$ whose remote pairs all lie in $T$ gives the constraint ``$\tau$ is not an induced $P_3$''; its scope consists of the owners of its soft pairs and the bits of its local hard pairs. For a triple $\tau$ let $P(\tau)$ be its non-remote pairs and $W(\tau)$ the union of $\mathcal M(u)$ over the ends $u$ of pairs of $P(\tau)$ (a set of members).
\begin{lemma}[Triple scopes]\label{lem:triplescope}
If $P(\tau)\ne\emptyset$, then $W(\tau)$ lies in a ball of radius $3$ of $I(H)$ around one member.
\end{lemma}
\begin{proof}
One pair $uv$: there are $A\in\mathcal M(u)$ and $B\in\mathcal M(v)$ equal or adjacent; every member of $\mathcal M(u)$ is within distance $1$ of $A$ (it shares $u$ with $A$), every member of $\mathcal M(v)$ within $1$ of $B$, so $W(\tau)$ lies within distance $2$ of $A$. Two pairs $uv,vw$ (they share $v$): there are $A\in\mathcal M(u)$, $B,B'\in\mathcal M(v)$, $C\in\mathcal M(w)$ with $d(A,B)\le1$, $d(B',C)\le1$, and $d(B,B')\le1$ as both contain $v$. Around $B$: $\mathcal M(v)$ is within $1$, $\mathcal M(u)$ within $2$, $C$ within $2$ and $\mathcal M(w)$ within $3$. Three pairs: two of them share a vertex and already involve all three vertices of $\tau$, so the argument for two pairs applies.
\end{proof}
Radius $3$ is attained: take six members forming a path $A'$--$A$--$B$--$B'$--$C$--$C'$ in $I(H)$ with $u\in A\cap A'$, $v\in B\cap B'$, $w\in C\cap C'$, so that $uv$ and $vw$ are local and $uw$ is remote; the minimal radius covering $W(\tau)$ is $3$.

\emph{Proof of Theorem~\ref{thm:L}.} Let $(\mathcal T,\{\beta_t\})$ be the given decomposition and $B_t$ the set of members within distance $3$ of $\beta_t$ in $I(H)$; $|B_t|\le(w+1)(1+\Delta+\Delta^2+\Delta^3)$ with $\Delta\le3(c-1)\le3w$. For a member $m$ the set $\{t:m\in B_t\}$ is the union of the subtrees $\{t:y\in\beta_t\}$ over the members $y$ within distance $3$ of $m$; these subtrees are chained along the connected ball, so the union is a subtree. A bit for a local pair $uv$ is placed in the bags $t$ with $\mathcal M(u)\cup\mathcal M(v)\subseteq B_t$, an intersection of subtrees, non-empty: take $t$ with $A\in\beta_t$ for $A\in\mathcal M(u)$ adjacent to some $B\in\mathcal M(v)$; then $\mathcal M(u)$ and $B$ are within $1$, $\mathcal M(v)$ within $2$. By Lemma~\ref{lem:triplescope}, all variables of a triple are in $B_t$ for any $t$ whose bag contains the centre member of the lemma. Hence this is a tree decomposition of the primal graph, with at most $|B_t|$ members and $\binom{3|B_t|}2$ bits per bag, and min-sum dynamic programming over it (at most $7^{|B_t|}2^{\binom{3|B_t|}{2}}$ states per bag) solves the CSP in time $f(w)n^{O(1)}$. Dropping triples with a remote pair outside $T$ is the definition of $F'(T)$. \hfill$\square$

(Checks: against an independent solver over pair variables on $280$ instances, no disagreement; the algorithm of Theorem~\ref{thm:B1} against exact $\OPT$ on $449$ (instance, budget) pairs, and, with an independent solver, on $237$ pairs, no mismatch; random tree-like instances needed radius $3$ and never $4$.)

\subsection{Proof of Theorem~\ref{thm:CB0}}\label{app:softCB}
Let $\ell=0$. Then $F'(T)>|H|$ iff there is no zero-excess assignment (states in $\{1,2,3\}$, no local hard edit) satisfying all triples whose remote pairs lie in $T$.
\begin{lemma}[Guards]\label{lem:guards}
Let $\tau$ be a triple with a remote pair, $\ell=0$. (i) If all three pairs are remote, or the non-remote pairs of $\tau$ are hard, then $\tau$ is a \emph{static test}: violated iff $\tau$ is an induced $P_3$ of $G$; it is active iff its (at most three) remote pairs lie in $T$. (ii) If exactly one pair $xz$ of $\tau$ is non-remote and it is soft with owner $M$, with remote pairs $xy,zy$, the triple is a \emph{guard}, a list on $M$ active iff $xy,zy\in T$: $t_{xz}(s_M)=1$ if both are edges, $=0$ if exactly one is an edge, no condition if none. For the three pairs of $M$ the lists are $\{2,3\},\{1,3\},\{3\}$ (for $t=1$) and $\{1\},\{2\},\{1,2\}$ (for $t=0$). (iii) If two pairs $uv,vw$ are non-remote they are not both soft (their owners would share $v$ and $uw$ would be a cross pair, hence local); one soft and one local hard gives a guard on the soft owner, active iff the remote pair $uw\in T$, of size $1$ or $2$; two local hard pairs give a static test.
\end{lemma}
\begin{proof}
Read off the definitions: a local hard pair is unedited at $\ell=0$ and a soft pair of a member in state $s$ is together iff $t_e(s)=1$; the triple is not a $P_3$ iff the number of edges is not $2$.
\end{proof}
Let $\Sigma_0$ be the set of zero-excess vectors that satisfy all triples with no remote pair (and are empty if one of these triples is a static violation). If $H$ has no soft triangle, $\Sigma_0$ is an intersection of unary lists and of the relations $R_{AB}$ (all four cross pairs are local), so it is closed under every majority operation (proof of Theorem~\ref{thm:sfA}).
\begin{lemma}[Helly lemma for lists of size at most two]\label{lem:sfH}
Let $|D|=3$ and $\Sigma\subseteq D^V$ closed under all majority operations. Let $L_i\subseteq D$, $|L_i|\le2$, for $i$ in a finite index set $I$, such that $\Sigma\cap\prod_iL_i=\emptyset$ but for every $i$ there is $x^i\in\Sigma$ with $x^i_l\in L_l$ for all $l\ne i$. Then $|I|\le3$.
\end{lemma}
\begin{proof}
Suppose $|I|\ge4$; then $x^i_i\notin L_i$. For $i,j,k\in I$ distinct and a majority operation $m$, $u=m(x^i,x^j,x^k)\in\Sigma$. For $l\notin\{i,j,k\}$ the three values lie in $L_l$, $|L_l|\le2$, so two coincide and $u_l\in L_l$. At $i$: if $x^j_i=x^k_i$ then $u_i=x^j_i\in L_i$. Otherwise $L_i=\{x^j_i,x^k_i\}$ and $(x^i_i,x^j_i,x^k_i)$ is a permutation $\pi_i$ of $D$; $u_i\in L_i$ iff $m(\pi_i)$ differs from the \emph{first} entry of $\pi_i$. Likewise at $j$ ($m(\pi_j)\ne$ second entry) and at $k$ (third entry). The values of $m$ on the six permutations are free, so a suitable $m$ exists unless $\pi_i=\pi_j=\pi_k$ (then $m(\pi)$ would have to avoid all three entries), i.e.\ unless $L_i,L_j,L_k$ are $2$-lists whose missing values are pairwise distinct. Among four coordinates two missing values coincide in every triple that contains them, so some triple has a suitable $m$; then $u\in\Sigma\cap\prod_iL_i$, a contradiction.
\end{proof}
\emph{Proof of Theorem~\ref{thm:CB0}.} Let $T$ be a minimal certificate. If one of its active static tests is violated, the violated triple has at most three remote pairs, all in $T$ by minimality, so $|T|\le3$. Otherwise the active guards of $T$ have empty intersection with $\Sigma_0$. Let $\mathcal G$ be a minimal family of active guards with $\Sigma_0\cap\bigcap\mathcal G=\emptyset$ and, for each member $i$, let $L_i$ be the intersection of the guards of $\mathcal G$ at $i$ ($|L_i|\le2$). Take a minimal set $I$ of coordinates with $\Sigma_0\cap\prod_{i\in I}L_i=\emptyset$; then for each $i\in I$ some $x^i$ exists as in Lemma~\ref{lem:sfH}, so $|I|\le3$, and by minimality of $\mathcal G$ all guards of $\mathcal G$ are at coordinates of $I$. At one coordinate, the partial intersections of its guards (non-redundant by minimality) strictly shrink: $\le2\to\le1\to0$. If some $L_i=\emptyset$, then $\mathcal G$ consists of the at most three guards at $i$; otherwise there are at most two guards at each coordinate. Hence $|\mathcal G|\le6$. Each guard has at most two remote pairs, and the remote pairs of $\mathcal G$ form a certificate contained in $T$, so equal to $T$: $|T|\le12$. \hfill$\square$

\subsection{The family \texorpdfstring{$G_N$}{G\_N} and the proof of Theorem~\ref{thm:CE}}\label{app:softCE}
\emph{Construction.} For $j=1,\dots,N$ take vertices $p_j,q_j,r_j,c_j,b_j$; add a vertex $a_1$ and put $a_j:=b_{j-1}$ for $j\ge2$. Members (centre in the middle)
\[A_j=(q_j,p_j,a_j),\qquad B_j=(q_j,r_j,b_j),\qquad C_j=(p_j,c_j,r_j).\]
The edges of $G_N$ are the pairs $q_jp_j,p_ja_j,q_jr_j,r_jb_j,p_jc_j,c_jr_j$ of the members and the hard edges $c_jq_j$ ($1\le j\le N$); all other pairs among these vertices are hard non-edges. Pendant vertices (in no member): $w_j$ adjacent to $p_j$ only ($1\le j\le N$), $w_{N+1}$ adjacent to $a_1$ only, $w_{N+2}$ adjacent to $b_N$ only; all their pairs are remote. Then $|V|=6N+3$ and $H_N$ has $3N$ members. Every vertex lies in at most two members ($p_j$ in $A_j,C_j$; $q_j$ in $A_j,B_j$; $r_j$ in $B_j,C_j$; $b_j$ in $B_j,A_{j+1}$), $I(H_N)$ is a chain of $N$ triangles joined by single edges ($\tw=2$), each triple $\{p_j,q_j,r_j\}$ is a soft triangle, and the hard edges form a matching, so no induced $P_3$ has three hard pairs: $H_N$ is pair-maximal (all by direct inspection; also checked by computer for $N\le3$).
\[T_N=\{w_jp_j,\,w_jr_j:1\le j\le N\}\cup\{w_{N+1}a_1,\,w_{N+1}p_1,\,w_{N+2}b_N,\,w_{N+2}r_N\},\qquad|T_N|=2N+4.\]

\begin{lemma}[Structure]\label{lem:GN}
Write $s_j=(A_j,B_j,C_j)$ for the states of block $j$. Then $\Sigma_0(G_N)$ (all zero-excess vectors satisfying every triple that has no remote pair) is the set of vectors with $s_j\in\{121,212,223\}$ for all $j$ and no $j<N$ with $s_j=212$ and $s_{j+1}=121$.
\end{lemma}
\emph{Proof (by computer, for the first blocks, and by isomorphism of the interior blocks).} A triple without remote pair lies in $W_j\cup W_{j+1}$ for some $j$, where $W_j$ is the vertex set of block $j$ (with $a_j,b_j$), by the definition of local. Inside a block the triples force the three states $121,212,223$ (the relation of a triangle block, by enumeration). The only triples with vertices of two blocks that carry soft pairs have both soft pairs at the glue vertex $b_j=a_{j+1}$ (the pairs $r_jb_j$ of $B_j$ and $p_{j+1}a_{j+1}$ of $A_{j+1}$) and the hard non-edge $r_jp_{j+1}$ as third pair: it is a $P_3$ iff both soft pairs stay edges, i.e.\ $B_j$ and $A_{j+1}$ are both in $\{1,3\}$, which within the allowed block states means $B_j=1$ ($s_j=212$) and $A_{j+1}=1$ ($s_{j+1}=121$). The other mixed triples have one soft pair and two non-edges. Enumeration over all $3^{3N}$ vectors for $N\le4$ gives $|\Sigma_0|=3,8,21,55$, equal to the transfer-matrix count, and contains every interior configuration (a block with both neighbours); interior blocks are isomorphic. \hfill$\square$

\emph{The active guards of $T_N$.} A triple is active iff its remote pairs lie in $T_N$. A triple containing $w_j$ ($j\le N$) has two remote pairs at $w_j$, so it is $\{w_j,p_j,r_j\}$; this is the guard of Lemma~\ref{lem:guards}(ii) with $xz=p_jr_j$, the pair $ac$ of $C_j$, $w_jp_j$ an edge and $w_jr_j$ a non-edge: $t_{ac}(C_j)=0$, i.e.\ $C_j\in\{1,2\}$. Likewise $\{w_{N+1},a_1,p_1\}$ gives $t_{bc}(A_1)=0$, i.e.\ $A_1=2$, and $\{w_{N+2},b_N,r_N\}$ gives $B_N=2$. No other triple is active, and each guard has exactly two remote pairs.

\emph{Proof of Theorem~\ref{thm:CE}.} \emph{Lower bound.} $F'(T_N)\ge|H|+1$: otherwise there would be a zero-excess vector in $\Sigma_0$ with $C_j\ne3$ for all $j$, $A_1=2$ and $B_N=2$. By Lemma~\ref{lem:GN} and $C_1\ne3$, block $1$ is $121$ or $212$; $A_1=2$ excludes $121$, so $s_1=212$, $B_1=1$. The link forbids $s_2=121$, and $C_2\ne3$ excludes $223$, so $s_2=212$; inductively $s_N=212$, $B_N=1$, contradicting $B_N=2$. \emph{Minimality.} Removing a pair $e\in T_N$ deactivates exactly one guard; for each guard a zero-excess vector exists: without the $A_1$-guard take all blocks $121$ (the link needs $s_j=212$, never); without the $B_N$-guard all blocks $212$ ($A_{j+1}=2$); without the $C_k$-guard, blocks $j<k$ equal to $212$, block $k$ equal to $223$ and blocks $j>k$ equal to $121$ (the link applies only after a block $212$, and the next block is then $212$ or $223$). These satisfy all remaining guards, and $F'(T_N\setminus\{e\})=|H|$. Also $F'(\emptyset)=|H|$ by $\Sigma_0\ne\emptyset$. \hfill$\square$

(The values for $N\le4$ were also recomputed by enumeration; the two further statements, $F'(T_N)=|H|+1$ for $N\le10$ and the connected family at $\ell=1$, are computer evidence, as stated in Appendix~\ref{sec:softcert}. $\OPT(G_N)=|H|+N$ was computed exactly for $N\le3$ by an independent subset dynamic program and for $N\le5$ by CP-SAT.)

\subsection{Proof of Lemma~\ref{lem:R5}}\label{app:softR5}
Identify $\{1,2\}$ with $\{0,1\}$ and put $R_5=\{(x,y,z)\in\{0,1\}^3:x+y+z\ne1\}$; then $R_k\cap\{1,2\}^3=R_5$ (the words $111,122,212,221,222$). Let $f$ be a conservative near-unanimity operation of arity $r\ge3$ preserving $R_k$ and all unary relations. Conservativity gives $f(\{1,2\}^r)\subseteq\{1,2\}$, so $f$ preserves $R_5$ on $\{0,1\}$. Let $\mathcal N_0=\{B\subseteq[r]:f(\mathbf 1_B)=0\}$ ($\mathbf 1_B$ the indicator vector). For $B,C\in\mathcal N_0$ apply $f$ to the three rows $\mathbf 1_{B\cup C},\mathbf 1_B,\mathbf 1_C$; every column is $(1,1,1),(1,1,0),(1,0,1)$ or $(0,0,0)$, hence in $R_5$, so $(f(\mathbf 1_{B\cup C}),0,0)\in R_5$, i.e.\ $f(\mathbf 1_{B\cup C})=0$ and $B\cup C\in\mathcal N_0$. Near-unanimity of arity $r\ge3$ gives $f(\mathbf 1_{\{i\}})=0$ (all but one argument equal $0$), so every singleton is in $\mathcal N_0$ and by union-closure $[r]\in\mathcal N_0$, i.e.\ $f(1,\dots,1)=0$, contradicting idempotence ($f(x,\dots,x)=x$). \hfill$\square$

The relation $R_k$ itself is realised as follows (by computer): six vertices $0,\dots,5$, members $A=(2,1,0)$, $B=(4,3,1)$, $C=(5,0,3)$ (centres $1,3,0$), hard non-edges $25,45$, and four further members that own the pairs $24,23,40,51$ (using fresh vertices; each of the vertices $0,\dots,4$ lies in three members, so $c=3$). The conjunction of the triples of the six vertices that contain none of these four pairs is $R_k$, with the states written $1,2,3$ for $ab,bc,ac$.

\subsection{Proof of Theorem~\ref{thm:B}}\label{app:softB}
\emph{Grid Tiling with $\le$}~\cite{Cygan15}: given sets $S_{ij}\subseteq[n]^2$ ($1\le i,j\le k$), choose $s_{ij}=(a_{ij},b_{ij})\in S_{ij}$ with $a_{ij}\le a_{i,j+1}$ and $b_{ij}\le b_{i+1,j}$; it is W[1]-hard in $k$ (and has no $f(k)n^{o(k)}$ algorithm under ETH). (This is the version of the cited chapter after exchanging the two coordinates.)

\emph{Construction.} For each cell $(i,j)$ create two \emph{arms}: $A_{ij}$ with one node per tile of $S_{ij}$ ordered by $(a,b)$, and $B_{ij}$ with one node per tile ordered by $(b,a)$. Row $i$ is the path obtained by concatenating $A_{i,k},A_{i,k-1},\dots,A_{i,1}$ (each in increasing order), column $j$ the path $B_{k,j},\dots,B_{1,j}$. The forest consists of these $2k$ paths. Obstruction paths: (1) every arm; (2) for every row $i$, $j<k$ and $\alpha\in[1,n-1]$ the set $\rho$ of nodes of $A_{i,j+1}$ with $a\ge\alpha+1$ together with the nodes of $A_{ij}$ with $a\le\alpha$ (a subpath of the row, as $A_{i,j+1}$ immediately precedes $A_{ij}$); (3) the same for columns with the $b$-coordinate and the arms $B$. Links: the node of the tile $e$ in $A_{ij}$ with the node of $e$ in $B_{ij}$ (a perfect matching). $L=k^2$.

\emph{Correctness.} ($\Rightarrow$) For a tiling use, for every cell, the link of its chosen tile: all arms are hit; a path $\rho$ of type (2) is hit iff $a_{ij}\le\alpha$ or $a_{i,j+1}\ge\alpha+1$, and $a_{ij}\le a_{i,j+1}$ gives this for every $\alpha$; columns likewise. ($\Leftarrow$) The $2k^2$ arms are pairwise disjoint and an operation meets at most two of them (a link only the two arms of its cell), so with $L=k^2$ every operation is a link, one per cell, hence a tile in every $S_{ij}$. If $a_{ij}>a_{i,j+1}$, the path $\rho$ for $\alpha=a_{i,j+1}$ contains neither chosen node (the chosen nodes of other cells are on other arms), a contradiction; columns likewise. An empty $S_{ij}$ gives an unhittable empty arm (NO, correctly), and an empty $\rho$ means $a_{ij}>a_{i,j+1}$ in every choice (NO). The size is polynomial and $L=k^2$, so an $f(L)n^{O(1)}$ algorithm would give an $f(k^2)n^{O(1)}$ algorithm for Grid Tiling with $\le$. \hfill$\square$

(Checks: own brute force and an independent solver, generic hitting-set branching on a smallest unhit path and naive enumeration of all $L$-subsets, on $1{,}710$ instances in total, including $k=2,3$ and $n=2,3,4$: no mismatch.)
